\cgeentry{cge_constraint}{Limits that bind in some scenarios and not in others}

\cgesec{Description}

\cmd{cge_constraint} puts an upper or a lower bound on elements of a model
variable, such as the imports of a commodity or the output of an activity, and
pairs each bound with an adjusting variable, which the closure fixes and which
may move in one direction only. In every solution one of two things holds:
either the limited variable is within its bound and the adjusting variable is
at its fixed level, or the limited variable is at its bound and the adjusting
variable has moved to keep it there. \manref{cge_solve} finds which, for each
constraint. \cmd{cge_constraint, list} displays the constraints, and
\cmd{cge_constraint, clear} removes them.

\cgesec{Quick start}

An import quota on manufactures, 2 per cent above base imports, with the
premium on imports adjusting
\begin{stcode}
cge_constraint QM, elements(com_man) upper(1.02) scale
\end{stcode}
A capacity limit on health services, 3 per cent above base output, with the
rent going to households
\begin{stcode}
cge_constraint QA, elements(act_hlt) upper(1.03) scale rents(hhd)
\end{stcode}
Under a fixed exchange rate, a limit of 30 on foreign savings, after which the
currency depreciates
\begin{stcode}
cge_constraint FSAV, upper(30) adjust(EXR) up
\end{stcode}
Government consumption is cut once government savings fall to zero
\begin{stcode}
cge_constraint GSAV, lower(0) adjust(GADJ) down
\end{stcode}
Display the constraints, with their regimes in the last solution; remove them
\begin{stcode}
cge_constraint, list
cge_constraint, clear
\end{stcode}

\cgesec{Syntax}

Add constraints
\begin{stsyntax}
\cmdb{cge_constraint} \argu{variable}\texttt{,}
\texttt{\{}\optarg{up}{per}{\#}\,|\,\optarg{lo}{wer}{\#}\texttt{\}}
[\argu{options}]
\end{stsyntax}
Display the constraints
\begin{stsyntax}
\cmdb{cge_constraint}\texttt{,} \opt{}{list}
\end{stsyntax}
Remove the constraints
\begin{stsyntax}
\cmdb{cge_constraint}\texttt{,} \opt{}{clear}
\end{stsyntax}

\begin{opttable}
\optgroup{Main}
\optarg{up}{per}{\#} & upper bound\\
\optarg{lo}{wer}{\#} & lower bound\\
\optarg{el}{ements}{list} & elements of \argu{variable} to bound; default every element\\
\opt{sc}{ale} & read the bound as a multiple of each element's base value\\
\optgroup{Adjusting variable}
\optarg{adj}{ust}{variable \textup{[}elements\textup{]}} & the variable that adjusts when the bound binds\\
\opt{}{up} & the adjusting variable may rise\\
\opt{}{down} & the adjusting variable may fall\\
\optarg{rent}{s}{account} & recipient of the rent; default the government\\
\end{opttable}

\argu{variable} is one model variable, in upper case. Exactly one of
\opt{up}{per()} and \opt{lo}{wer()} is required, and at most one of
\opt{}{up} and \opt{}{down} is allowed.

\cgesec{Options}

\begin{optdesc}
\item[\optarg{up}{per}{\#} \textrm{or} \optarg{lo}{wer}{\#}] gives the bound,
as a level or, with \opt{sc}{ale}, as a multiple of the base value.

\item[\optarg{el}{ements}{list}] names the elements of the limited variable;
there is one constraint for each. Without it, every existing element of a
variable with one index is bounded. A variable with two indexes takes one
element of each, as in \cmd{elements(lab act_agr)}, and a variable without
index takes none.

\item[\opt{sc}{ale}] reads the bound as a multiple of each element's base
value. An element whose base value is zero needs a bound given as a level.

\item[\optarg{adj}{ust}{variable \textup{[}elements\textup{]}}] names the
adjusting variable, which must be fixed under the closure. Without elements it
takes the element of the limited variable, when the two variables have the
same index; a variable without index, or an adjusting variable whose elements
are given, serves one constraint. A bound on \cmd{QM}, \cmd{QE} or \cmd{QA} is
paired by default with the rent of the same element: \cmd{RM}, \cmd{RE} or
\cmd{RA}. Other limited variables need \opt{adj}{ust()}.

\item[\opt{}{up} \textrm{or} \opt{}{down}] gives the direction in which the
adjusting variable may move from its fixed level. Rents rise by default; for
other adjusting variables the direction is required.

\item[\optarg{rent}{s}{account}] names the recipient of the rent of
\cmd{RM}, \cmd{RE} or \cmd{RA}: a factor, a household or enterprise, the
government or the rest of the world, by the label of its account. The default
is the government.

\item[\opt{}{list}] displays the constraints, with their regime in the last
solution: \cmd{binding}, \cmd{not binding}, or \cmd{not solved} for a
constraint added since.

\item[\opt{}{clear}] removes every constraint and sends every rent to the
government.
\end{optdesc}

\cgesec{Remarks and examples}

A model with only equalities cannot represent a limit that binds in some
scenarios and not in others: an import quota, the capacity of an activity, a
ceiling on borrowing or a floor under a wage. \cmd{cge_constraint} writes each
limit as a complementarity condition between the limited variable and an
adjusting variable, and \manref{cge_solve} solves the conditions with the
rest of the model, so that the solver finds which limits bind; none is chosen
in advance.

Each constraint adds one equation and frees one variable, the adjusting
variable, so the system stays square. The limited variable must be free and
the adjusting variable fixed under the closure, and every later closure set by
\manref{cge_closure} is checked against the constraints. A constraint on the
same element with the same type of bound replaces the earlier one, and each
element of an adjusting variable serves one constraint. Constraints are kept
by \cmd{cge_shock, reset}, in dynamic runs, in sensitivity analysis and by
\manref{cge_save}; they are written into the GAMS program by
\manref{cge_gams}, which solves them as a mixed complementarity problem.
\manref{cge_calibrate} removes them.

\cgesubsec{Rents}

Quotas and capacity limits work through rents, which are zero in the base and
fixed by every closure unless a constraint frees them:

\begin{cgetable}{p{0.12\linewidth}p{0.82\linewidth}}
\toprule
Variable & Description\\
\midrule\endhead
\bottomrule\endlastfoot
$\V{RM}_c$ & premium on imports of commodity $c$, per unit of the world price\\
$\V{RE}_c$ & premium on exports of commodity $c$, per unit of the world price\\
$\V{RA}_a$ & rent on the output of activity $a$, per unit of its price\\
\end{cgetable}

A premium is a wedge between the world price and the domestic price, like a
tariff, but its revenue is a rent, paid to the recipient given in
\opt{rent}{s()}: the government by default, as with auctioned quota licences;
a household or enterprise; a factor, such as the owners of capital; or the
rest of the world, as with a voluntary export restraint, whose rents foreign
exporters keep. A rent set by hand with \manref{cge_shock}, without a
constraint, acts as a tax whose revenue goes to the recipient: for example
\cmd{cge_shock RM, elements(com_man) value(0.05)} acts as a tariff of 5 per
cent.

\cgesubsec{Typical constraints}

\begin{cgetable}{p{0.22\linewidth}llp{0.44\linewidth}}
\toprule
Limit & Limited & Adjusting & Effect when it binds\\
\midrule\endhead
\bottomrule\endlastfoot
import quota & $\V{QM}_c \le b$ & $\V{RM}_c$, up & the premium rises until imports equal the quota\\
export quota & $\V{QE}_c \le b$ & $\V{RE}_c$, up & the premium rises until exports equal the quota\\
capacity & $\V{QA}_a \le b$ & $\V{RA}_a$, up & the activity price rises above cost; the rent goes to the owners of the capacity\\
borrowing limit under a peg & $\V{FSAV} \le b$ & $\V{EXR}$, up & the peg holds until foreign savings reach the limit; the currency then depreciates\\
deficit ceiling & $\V{GSAV} \ge b$ & $\V{GADJ}$, down & government consumption is cut once the deficit reaches the ceiling\\
deficit ceiling, taxes & $\V{GSAV} \ge b$ & $\V{TINSADJ}$, up & direct tax rates rise instead\\
nominal wage floor & $\V{W}_f \ge b$ & $\V{QFS}_f$, down & employment falls once the wage reaches the floor\\
\end{cgetable}

The commands are, for example,
\begin{stcode}
cge_constraint QM, elements(com_man) upper(1.02) scale
cge_constraint QA, elements(act_hlt) upper(1.03) scale rents(hhd)
cge_constraint FSAV, upper(30) adjust(EXR) up
cge_constraint GSAV, lower(0) adjust(GADJ) down
cge_constraint W, elements(lab) lower(1) scale adjust(QFS) down
\end{stcode}
The borrowing limit needs the closure \cmd{external(fixexr)}, under which
$\V{EXR}$ is fixed and $\V{FSAV}$ free, and the wage floor a closure under
which the wage is free and the supply of labour fixed, such as \cmd{mobile}.

\cgesubsec{Example 1: a capacity limit}

The health layout of \manref{cge_synthetic} is mapped, with its social health
insurance fund, and calibrated as in the examples of \manref{cge_solve}.
Government purchases of health services rise by 20 per cent, first without and
then with a limit on what health providers can deliver, 3 per cent above the
base:
\stlog{cons_cap}
A new constraint is \cmd{not solved} until the model is solved. After a solve,
\manref{cge_solve} displays the regime of each constraint and the values of
its two variables: the limit binds, output is at the bound of 30.9, and the
rent on output is 5.9 per cent of the activity price. The report shows the
effect of the limit:
\stlog{cons_cap_res}
Without the limit, the output of health services rises by 6.6 per cent and
their price by 0.2 per cent. With it, output rises by exactly 3 per cent and
the price by 6 per cent: the rent on output fills the gap between price and
cost, and goes to households, the owners of the capacity here.

\cgesubsec{Example 2: a ceiling on the deficit}

The same spending, with a government that may not let its savings fall below
their base value: spending on other goods and services is cut once the limit
is reached.
\stlog{cons_deficit}
The command echo shows the local macro as typed; its value, 1, is the base
value of government savings. The constraint binds, and the scaling factor of
government consumption, $\V{GADJ}$, falls by 4.6 per cent.

\cgesubsec{Example 3: a quota that binds in one scenario and not in another}

An import quota on manufactures, 2 per cent above base imports, does not bind
in the base. When import tariffs are removed, demand for imports rises and the
quota binds:
\stlog{cons_quota}
Imports are at the quota, and the premium on imports is 4.1 per cent of the
world price; the rent goes to the government:
\stlog{cons_quota_res}

\cgesubsec{Example 4: a borrowing limit under a peg}

Foreign savings are fixed under the default closure and cannot be bounded. Under
a fixed exchange rate they are free, and the peg can hold until they reach a
limit, here 5 per cent above the base, with the tariff removal of Example~3
still in force:
\stlog{cons_peg}
The limit binds: foreign savings are at 14.7, and the exchange rate has
depreciated by 4.7 per cent. A closure that makes the constraint invalid is
itself rejected: with this constraint, \cmd{cge_closure, external(flexexr)}
fails, since it would fix foreign savings and free the exchange rate.

\cgesubsec{The minimum wage}

The minimum real wage of the closure \cmd{minwage} of \manref{cge_closure} is
a complementarity condition of the same kind, with its own calibration and its
own smoothed condition (\secref{sec:minwage}). It is not set with
\cmd{cge_constraint}.

\cgesec{Stored results}

\cmd{cge_constraint} stores the following in \cmd{r()}:
\begin{storedres}
\resgroup{Scalars}
\res{r(N_constraints)} & number of constraints\\
\res{r(N_binding)} & number binding in the last solution (\opt{}{list} only)\\
\end{storedres}
\manref{cge_solve} stores the regimes of a solution in \res{r(constraints)},
with \res{r(N_constraints)} and \res{r(N_binding)}.

\cgesec{Methods and formulas}

A constraint $k$ pairs a limited variable $Y_k$, one element of a model
variable, and its bound $b_k$ with an adjusting variable $X_k$, which the
closure fixes at $\bar{x}_k$, its current value including any shock. With
\opt{sc}{ale}, $b_k$ is the number given times the base value $Y^0_k$. The
conditions are
\begin{equation*}
g_k = \sigma_k\,\frac{b_k - Y_k}{s^Y_k} \ge 0, \qquad
h_k = \delta_k\,\frac{X_k - \bar{x}_k}{s^X_k} \ge 0, \qquad
g_k\,h_k = 0,
\end{equation*}
with $\sigma_k = 1$ for an upper bound and $-1$ for a lower bound,
$\delta_k = 1$ for \opt{}{up} and $-1$ for \opt{}{down}, and $s$ the base
magnitudes of the variables (one where the base value is zero).

\manref{cge_solve} solves them in two stages. First each condition is written
in the smoothed Fischer--Burmeister form
\begin{equation*}
g_k + h_k - \sqrt{g_k^2 + h_k^2 + 2\mu} = 0, \qquad \mu = 10^{-8},
\end{equation*}
and solved with the rest of the model. Then each constraint takes the regime
of that solution, binding ($g_k = 0$) if $h_k > g_k$ and otherwise not binding
($h_k = 0$), and the model is solved again with these exact equations. A
constraint whose other condition then fails ($h_k < 0$ when binding,
$g_k < 0$ when not) changes regime, and the model is solved again, for at most
ten rounds.

The rent on each element,
\begin{equation*}
R^M_c = \V{RM}_c\,\V{pwm}_c\,\V{EXR}\,\V{QM}_c, \qquad
R^E_c = \V{RE}_c\,\V{PWE}_c\,\V{EXR}\,\V{QE}_c, \qquad
R^A_a = \V{RA}_a\,\V{PA}_a\,\V{QA}_a,
\end{equation*}
is income of its recipient, and rents paid abroad are payments to the rest of
the world. The premiums enter the import and export prices, \eqref{eq:P1} and
\eqref{eq:P2}, and the rent on output the activity price, \eqref{eq:P7}.
\Chapref{chap:constraints} gives the details (\secref{sec:rents},
\secref{sec:cons-def} and \secref{sec:cons-solution}).

\cgesec{Also see}

\begin{alsosee}
\seeentry{cge_solve}{Solve the model}
\seeentry{cge_closure}{Choose the model closure}
\seeentry{cge_shock}{Change a parameter or a fixed variable}
\seeentry{cge_gams}{Write the model as a GAMS program, and solve it in GAMS}
\seechapter{chap:constraints}{Bounds and complementarity}
\end{alsosee}
